\begin{equation}
    \begin{array}{lrll}
        \text{Formulas ($o$):}  & \phi, \psi & ::= X^o                   \quad  & \text{(Variabili proposizionali)} \\
                                &            & \mid \dots                \quad  & \text{(Connettivi IPC)}           \\
                                &            & \mid (\forall X^o \phi)^o \quad  & \text{(Binder Universale)}        \\
                                &            & \mid (\exists X^o \phi)^o \quad  & \text{(Binder Esistenziale)}
    \end{array}
    \tag{$\text{IPC}^2$}
    \label{eq:ipc2_grammar}
\end{equation}

\paragraph{Regole di formazione.}
Le regole di buona formazione del secondo ordine estendono il concetto di sensatezza strutturale abilitando l'impredicatività 
logica entro binari controllati \cite{hindley1997}. Il binder puo vincolare variabili associate a valori di verità ($o$). L'istanza 
di buona formazione garantisce che una formula possa quantificare validamente sopra l'intero insieme dei tipi proposizionali (incluso 
il tipo della formula definenda stessa) senza generare le contraddizioni autoreferenziali, poiché la stratificazione sintattica isola 
il metalinguaggio dal linguaggio oggetto.